\section*{الفصل \ref{ch:b1:optical-mineralogy} --- علم المعادن البصري}

\begin{solution}{exo:b1:optical-mineralogy:1}
الكوارتز: $\qty{30000}{nm}\times0.009 = \qty{270}{nm}$، وهو أبيض من الرتبة الأولى.
الكالسيت: $\qty{30000}{nm}\times0.172 = \qty{5160}{nm}$، أي بعيدًا خارج المخطط:
أبيض باهت من الرتبة العالية، يميل قليلًا إلى الوردي أو الأخضر.
\end{solution}

\begin{solution}{exo:b1:optical-mineralogy:2}
$I/I_{\max} = \sin^2 2\theta$: $0$ عند \ang{0}؛ $\sin^2\ang{45} = 0.5$ عند \ang{22.5}؛ $1$ عند \ang{45}؛ $0.5$ عند \ang{67.5}؛ $0$ عند \ang{90}.
\end{solution}

\begin{solution}{exo:b1:optical-mineralogy:3}
\omterm{def:b1:rock-forming-minerals:groups}{الجارنيت} مكعّب، ومن ثم \omterm{def:b1:optical-mineralogy:anisotropic}{متساوي الخواص}: يخرج الضوء منه مهتزًا كما دخل، على طول المستقطِب، فيوقفه المحلِّل؛ و$\Delta = 0$ في كل وضع. وحبيبة كوارتز مقطوعة عموديًا على محورها البصري لها أيضًا $n_1 = n_2$ للضوء المنتشر على طول ذلك المحور، ف$\Delta = 0$: فهي أيضًا تبقى معتمة. وفي \omterm{def:b1:field-lab-methods-1:microscope}{الضوء المستقطَب المستوي} يكون \omterm{def:b1:rock-forming-minerals:groups}{للجارنيت} بروز مرتفع جدًا ($1.83$ مقابل $1.54$)، والكوارتز يكاد لا بروز له.
\end{solution}

\begin{solution}{exo:b1:optical-mineralogy:4}
$|n - 1.54|$: الألماندين $0.29$، والفورستيريت $0.11$، والأورثوكليز $0.02$، والكوارتز $0.01$؛ \omterm{def:b1:optical-mineralogy:relief}{فالبروز} بهذا الترتيب، من الأعلى إلى الأدنى. للأورثوكليز معامل أدنى: بروز سالب، \omterm{def:b1:optical-mineralogy:relief}{وخط بيك} يخرج من الحبيبة إلى وسط التثبيت عند رفع البؤرة.
\end{solution}

\begin{solution}{exo:b1:optical-mineralogy:5}
$\delta = 650/30000 = 0.022$. \omterm{def:b1:rock-forming-minerals:mineral}{معدن} أخضر \omterm{def:b1:optical-mineralogy:pleochroism}{متعدد الألوان}، له انشقاقان عند \ang{56} و\ang{124}، وانكسار مزدوج قرب $0.020$: إنه الهورنبلند.
\end{solution}

\begin{solution}{exo:b1:optical-mineralogy:6}
$\tan\ang{30} = 0.5774$. الكوارتز: $(0.013/1.54)\times0.5774 = \qty{0.00487}{rad}
= \ang{0.28}$. الفورستيريت: $(0.111/1.54)\times0.5774 = \qty{0.0416}{rad} =
\ang{2.38}$. النسبة $0.111/0.013 = 8.5$. وبالضبط: $\sin r = 1.54\sin\ang{30}/1.651
= 0.4664$، $r = \ang{27.80}$، فيكون الانحراف \ang{2.20}: الصيغة من \omterm{def:b1:optical-mineralogy:interference-colour}{الرتبة} الأولى تبالغ بنحو \qty{8}{\percent} لفرق في المعاملات كبير بهذا القدر.
\end{solution}

\begin{solution}{exo:b1:optical-mineralogy:7}
$270 + 540 = \qty{810}{nm}$، وهو أخضر من الرتبة الثانية، حين تتوازى الاتجاهات البطيئة للحبيبة والصفيحة (فتتجمع التأخيرات)؛ و$|270 - 540| =
\qty{270}{nm}$، أبيض من \omterm{def:b1:optical-mineralogy:interference-colour}{الرتبة} الأولى مرة أخرى، حين يتوازى الاتجاه البطيء للحبيبة مع الاتجاه السريع للصفيحة.
\end{solution}

\begin{solution}{exo:b1:optical-mineralogy:8}
عند السمك الكامل $\Delta = \qty{30000}{nm}\times0.052 = \qty{1560}{nm}$. ومن الحافة ($\Delta = 0$) نحو الداخل يمرّ \omterm{def:b1:optical-mineralogy:retardation}{التأخير} ب\num{542} و\qty{1074}{nm} دون \qty{1605}{nm}: هدبتان أرجوانيتان محمّرتان.
\end{solution}

\begin{solution}{exo:b1:optical-mineralogy:9}
\omterm{def:b1:optical-mineralogy:extinction}{زاوية الانطفاء} المقاسة على حبيبة تتوقف على كيفية قطعها: لا تُظهر الزاوية الأكبر إلا المقاطع ذات توجّه معيّن، وتُظهر غيرها أقل، حتى الصفر. وأكبر قيمة، نحو \ang{37}، هي التي تُقارن: إنها تطابق الديوبسيد (نحو \ang{38})، لا الهورنبلند (في أكبر الأحوال \ang{19}).
\end{solution}

\begin{solution}{exo:b1:optical-mineralogy:10}
مع محلِّل على طول $x$، تكون مركبتا $x$ للموجتين
$E_0\cos\theta\cdot\cos\theta$ و$-E_0\sin\theta\,e^{-i\varphi}\cdot(-\sin\theta)$،
فيكون $A = E_0(\cos^2\theta + \sin^2\theta\,e^{-i\varphi})$ و
\[
  |A|^2 = E_0^2\bigl(\cos^4\theta + \sin^4\theta + 2\sin^2\theta\cos^2\theta\cos\varphi\bigr)
  = E_0^2\bigl(1 - 2\sin^2\theta\cos^2\theta\,(1 - \cos\varphi)\bigr),
\]
مستعملين $\cos^4\theta + \sin^4\theta = 1 - 2\sin^2\theta\cos^2\theta$. وبما أن
$2\sin^2\theta\cos^2\theta = \frac12\sin^2 2\theta$ و$1 - \cos\varphi =
2\sin^2(\varphi/2)$، فإن $I_\parallel = I_0(1 - \sin^2 2\theta\sin^2(\pi\Delta/\lambda))$.
عند كل طول موجي، $I_\parallel + I_\perp = I_0$: فالألوان بين مستقطِبين متوازيين هي متممات الألوان بين \omterm{def:b1:field-lab-methods-1:microscope}{المستقطِبين المتقاطعين}. وحبيبة لها \qty{550}{nm}، أرجوانية بين \omterm{def:b1:field-lab-methods-1:microscope}{المستقطِبين المتقاطعين}، تكون مخضرّة مصفرّة باهتة بين المستقطِبين المتوازيين.
\end{solution}

\begin{solution}{exo:b1:optical-mineralogy:11}
(أ) $\Delta/\lambda$ يمتد من $5000/700 = 7.14$ إلى $5000/400 = 12.5$؛ فيمرّ بالأعداد الصحيحة من 8 إلى 12: خمسة أصفار. (ب) الأصفار عند $\lambda_k = \Delta/k$؛ وصفران متتاليان يبعدان $\Delta/k - \Delta/(k + 1) =
\Delta/(k(k + 1)) \approx \lambda^2/\Delta$، أي $550^2/5000 =
\qty{60.5}{nm}$ قرب \qty{550}{nm}. (ج) خمسة نطاقات داكنة ضيقة منتشرة عبر الطيف تُزيل قليلًا من كل لون، ويمرّ الباقي بوزن متوسط نصف: فتدرك العين لونًا باهتًا شبه محايد.
\end{solution}

\begin{solution}{exo:b1:optical-mineralogy:12}
متوسط $\sin^2 2\theta$ على دورة كاملة هو $\frac1{2\pi}\int_0^{2\pi}\sin^2 2\theta\,
\mathrm d\theta = \frac12$، فيكون متوسط الشدة
$\frac12 I_0\sin^2(\pi\Delta/\lambda)$. وهي تحسب $\delta = \Delta/30000$
بدلًا من $\Delta/25000$ (بالنانومتر): فكل قيمها هي $25/30 = 0.83$ من القيم الحقيقية، أي أدنى منها بنحو \qty{17}{\percent}.
\end{solution}

\begin{solution}{pb:b1:optical-mineralogy:1}
\textbf{1.} \omterm{def:b1:rock-forming-minerals:mineral}{المعدن} B: الحواف الداكنة تعني بروزًا مرتفعًا؛ وحدود A بالكاد مرئية تعني بروزًا منخفضًا.

\textbf{2.} A: $(1.544 + 1.553)/2 = 1.5485$، $|n - n_m| = 0.0085$. أما B:
$(1.664 + 1.672 + 1.694)/3 = 1.6767$، $|n - n_m| = 0.1367$.

\textbf{3.} الانحرافات متناسبة مع $|n - n_m|$: $0.1367/0.0085 =
16.1$.

\textbf{4.} معامل انكسار B أعلى من وسط التثبيت: بروز موجب.

\textbf{5.} \omterm{def:b1:optical-mineralogy:pleochroism}{تعدّد الألوان}؛ البيوتيت.

\textbf{6.} انشقاقان عند \ang{56} و\ang{124}: \omterm{def:b1:rock-forming-minerals:groups}{أمفيبول}.

\textbf{7.} $\qty{30000}{nm}\times0.009 = \qty{270}{nm}$: أبيض من الرتبة الأولى.

\textbf{8.} تُقطع كل حبيبة بتوجّهها الخاص وتُظهر قيمة $|n_1 - n_2|$ بين 0 و$\delta$؛ ولا تُظهر $\delta$ نفسها إلا الحبيبات المقطوعة موازية للمحور البصري. وأعلى لون بين حبيبات كثيرة هو الذي يُستعمل.

\textbf{9.} $d = 342/0.009 = \qty{38000}{nm} = \qty{38}{\micro m}$.

\textbf{10.} $(38 - 30)/30 = \qty{26.7}{\percent}$.

\textbf{11.} $\sin^2(\pi\times342/550) = 0.861$؛ وعند \qty{30}{\micro m}،
$\sin^2(\pi\times270/550) = 0.999$.

\textbf{12.} من $\qty{38000}{nm}\times0.008 = \qty{304}{nm}$ إلى
$\qty{38000}{nm}\times0.013 = \qty{494}{nm}$: من الأبيض إلى الأحمر البرتقالي، كلها دون نهاية \omterm{def:b1:optical-mineralogy:interference-colour}{الرتبة} الأولى عند \qty{542}{nm}.

\textbf{13.} بين 1074 و\qty{1605}{nm}: \omterm{def:b1:optical-mineralogy:interference-colour}{الرتبة} الثالثة.

\textbf{14.} $\delta = 1140/38000 = 0.030$.

\textbf{15.} $1140/30000 = 0.038$، وهو ما كان سيوحي \omterm{def:b1:rock-forming-minerals:groups}{بالأوليفين}. \omterm{def:b1:rock-forming-minerals:groups}{فالأوليفين} المغنيسيومي والكوارتز يتفاعلان ليعطيا \omterm{def:b1:rock-forming-minerals:groups}{البيروكسين}،
\ce{Mg2SiO4 + SiO2 -> 2 MgSiO3}؛ وصخرة تبلورت مع الكوارتز لا ينبغي أن تُبقي عليهما معًا.

\textbf{16.} $0.030$ هو \omterm{def:b1:optical-mineralogy:birefringence}{الانكسار المزدوج} للديوبسيد؛ وبروزه المرتفع، وغياب لونه، وانشقاقاه القريبان من \ang{87} كلها تتفق: B هو الكلينوبيروكسين، أي الديوبسيد.

\textbf{17.} $\qty{30000}{nm}\times0.030 = \qty{900}{nm}$: أصفر من الرتبة الثانية.

\textbf{18.} بمقدار \qty{8}{\micro m}؛ فلا يُظهر الكوارتز عندئذ أكثر من \qty{270}{nm}،
أبيض من الرتبة الأولى.

\textbf{19.} \omterm{def:b1:optical-mineralogy:twin}{التوأمة المتعددة التركيب}: D هو \omterm{def:b1:rock-forming-minerals:groups}{البلاجيوكلاز}.

\textbf{20.} تركيب \omterm{def:b1:rock-forming-minerals:groups}{البلاجيوكلاز} (محتواه من الأنورثيت).

\textbf{21.} \omterm{def:b1:rock-forming-minerals:groups}{أمفيبول} زاوية انطفائه \ang{17}: هورنبلند
(\ang{15}--\ang{19})، لا \omterm{def:b1:rock-forming-minerals:groups}{بيروكسين} (\ang{38} للديوبسيد).

\textbf{22.} $\qty{38000}{nm}\times0.020 = \qty{760}{nm}$: أخضر من
الرتبة الثانية.

\textbf{23.} هي مقطوعة عموديًا على \omterm{def:b1:optical-mineralogy:optic-axis}{المحور البصري}: معاملاها متساويان للضوء المنتشر على طول المحور، $\Delta = 0$، فتبقى معتمة كحبيبة متساوية الخواص.

\textbf{24.} A كوارتز، وB ديوبسيد، وC بيوتيت، وD \omterm{def:b1:rock-forming-minerals:groups}{بلاجيوكلاز}، وE هورنبلند: صخر ناري متوسط يحوي بعض الكوارتز، مثل الديوريت الكوارتزي.

\textbf{25.} سُمك المقطع \textbf{\qty[text-series-to-math=true]{38}{\micro m}}.
\end{solution}
